\section*{ВВЕДЕНИЕ}
\addcontentsline{toc}{section}{ВВЕДЕНИЕ}

Корпоративные коммуникации -- это не просто переписка, это кровеносная система бизнеса, требующая надёжности и эффективности. Закрытый корпоративный мессенджер - это специализированное веб-приложение, созданное для решения именно этих задач. 

Ключевые преимущества решения:
\begin{itemize}
	\item система аутентификации (логин/пароль);
	\item ролевое управление группами (модератор/владелец/участник);
	\item шифрование передаваемых данных (HTTPS);
	\item поддержка многопользовательских чатов;
	\item история сообщений с временными метками;
	\item отправка файлов всех форматов.
\end{itemize}

\emph{Цель работы} — разработка корпоративного мессенджера, полностью независимого от санкций и обеспечивающего бесперебойную внутреннюю коммуникацию. Безопасную платформу с функциями:
\begin{itemize}
	\item общий и групповые чаты;
	\item приватная переписка;
	\item управление ролями в группах;
	\item поиск по истории сообщений;
	\item отправка файлов;
	\item регистрация и авторизация.
\end{itemize}

\textit{Структура и объем работы}. Отчет состоит из введения, 4 разделов основной части, заключения, списка использованных источников, 2 приложений. Текст выпускной квалификационной работы равен \formbytotal{lastpage}{страниц}{е}{ам}{ам}.

\textit{Во введении} сформулирована цель работы, поставлены задачи разработки, описана структура работы, приведено краткое содержание каждого из разделов.

\textit{В первом разделе} на стадии описания анализа предметной области приводится сбор информации о возможностях корпоративных мессенджеров, а также производится анализ существующих аналогов.

\textit{Во втором разделе} на стадии технического задания приводятся требования к функционалу разрабатываемого продукта.

\textit{В третьем разделе} на стадии технического проектирования представлены проектные решения, и их обоснования.

\textit{В четвертом разделе} приводится список классов и их методов, использованных при разработке программы, производится тестирование разработанной системы.

\textit{В заключении} излагаются основные результаты работы, полученные в ходе разработки.

В приложении А представлен графический материал. В приложении Б представлены фрагменты исходного кода.